% !TEX program = xelatex
\documentclass[12pt,a4paper]{article}

\usepackage{fontspec}
\usepackage[spanish,es-tabla]{babel}
\usepackage{geometry}
\usepackage{setspace}
\usepackage{microtype}

\geometry{
	left=3cm,
	right=3cm,
	top=3cm,
	bottom=3cm
}

\setmainfont{Times New Roman}
\onehalfspacing
\setlength{\parindent}{0pt}
\setlength{\parskip}{0.8em}

\begin{document}

\begin{center}
	{\Large\bfseries Contribución de la investigación}
\end{center}

\vspace{1em}

La contribución de esta investigación no consiste en proponer una teoría completamente nueva sobre la relación entre recursos naturales y desarrollo, ni en afirmar que cada variable utilizada sea inédita en la literatura. Existe ya una línea empírica reciente que examina vínculos muy próximos entre recursos naturales, rentas extractivas, dependencia de recursos y complejidad económica, tanto en muestras regionales como en paneles amplios de países (Ajide, 2022; Avom et al., 2022; Canh et al., 2020; Destek et al., 2023; Kelly et al., 2024; Owjimehr \& Jamshidi, 2024; Zhang et al., 2023). En diálogo con esos trabajos, el aporte de esta tesis radica en reorganizar empíricamente el problema de la transformación estructural en economías dependientes de recursos naturales no renovables, distinguiendo entre el criterio que define el universo de análisis, los mecanismos explicativos y los resultados que se busca explicar. En ese sentido, el trabajo no se presenta como una réplica ni como una sustitución de dicha literatura, sino como una aplicación comparada orientada a identificar qué canales se asocian con trayectorias distintas de complejidad y diversificación exportadora dentro de economías con una inserción extractiva común.

El primer aporte consiste en separar la dependencia externa de recursos naturales de la intensidad macroeconómica de las rentas extractivas. Varios estudios analizan la relación entre rentas, abundancia o agotamiento de recursos naturales y complejidad económica en muestras amplias de países, países en desarrollo o economías ricas en recursos, y algunos distinguen incluso entre petróleo, minería y gas (Canh et al., 2020; Hoang et al., 2023; Kelly et al., 2024; Li et al., 2024; Ul-Durar et al., 2023; Valentine et al., 2024; Zhang et al., 2023). Esta tesis toma ese punto de partida, pero utiliza la dependencia externa, medida por la participación de las exportaciones de recursos naturales no renovables en la canasta exportadora, como criterio para delimitar la muestra. En cambio, las rentas de recursos naturales como porcentaje del PIB se incorporan como una variable explicativa dentro del modelo econométrico. Esta distinción permite evitar una confusión frecuente entre la dependencia de recursos como condición estructural de partida y las rentas extractivas como mecanismo económico que puede afectar las trayectorias de transformación estructural.

El segundo aporte se encuentra en las variables explicadas. La literatura más cercana se concentra principalmente en el Índice de Complejidad Económica (ECI) como resultado para estudiar si las rentas, la abundancia o la dependencia de recursos naturales limitan la sofisticación productiva de las economías (Ajide, 2022; Canh et al., 2020; Hoang et al., 2023; Owjimehr \& Jamshidi, 2024; Valentine et al., 2024; Zhang et al., 2023). Además, las aplicaciones recientes del ECI muestran que la complejidad se relaciona con dimensiones estructurales como energía, innovación, productividad, dependencia de recursos y capacidades mineras (Destek et al., 2023; Kazemzadeh et al., 2023; Sweet \& Eterovic, 2019; Valverde Carbonell et al., 2023). En este marco, la tesis aproxima la transformación estructural desde la inserción externa. La variable dependiente principal es el ECI, entendido como una medida de sofisticación exportadora y de capacidades productivas reveladas. De manera complementaria, se incorpora la diversificación exportadora, definida como $DIVX=1-HHI$, donde $HHI$ corresponde al índice de concentración de Herfindahl--Hirschman. Esta segunda variable no reemplaza al ECI, sino que permite observar si los mecanismos asociados a la complejidad económica también se relacionan con una canasta exportadora menos concentrada.

El tercer aporte consiste en articular en un mismo marco empírico distintos canales explicativos que la literatura cercana suele tratar de manera separada. Algunos estudios enfatizan el papel de las instituciones, la estabilidad política o la calidad institucional en economías ricas en recursos o en paneles amplios de países (Ajide, 2022; Avom et al., 2022; Avom \& Ndoya, 2024; Olaniyi \& Odhiambo, 2025a, 2025c; Owjimehr \& Jamshidi, 2024); otros destacan el desarrollo financiero, la innovación, el capital humano, los derechos de patente o el nivel de ingreso como condiciones que pueden modificar la relación entre recursos y complejidad (Alvarado et al., 2023; Li et al., 2024; Mayer et al., 2023; Sweet \& Eterovic, 2019; Valentine et al., 2024; Yalta \& Yalta, 2021). A partir de estos antecedentes, el modelo incorpora un canal institucional, asociado a la calidad institucional y a su interacción con las rentas extractivas; un canal de abundancia, vinculado con petróleo, gas y minería; un canal estructural, relacionado con concentración y especialización exportadora; un canal macroeconómico, asociado a volatilidad y tipo de cambio real; un canal de capacidades productivas, aproximado mediante capital humano, innovación y conectividad; y canales fiscales y financieros. El aporte no reside en afirmar que todos estos canales sean novedosos por separado, sino en evaluarlos de forma conjunta para comparar cómo se relacionan con resultados distintos de complejidad y diversificación en economías dependientes de recursos no renovables.

Finalmente, la tesis aporta una estrategia empírica comparada que busca diferenciarse de los estudios existentes no por el simple uso de datos de panel, sino por la forma en que define y organiza el universo de análisis. En efecto, ya existen aplicaciones de panel sobre recursos naturales y complejidad económica, incluyendo estudios con muestras globales, países en desarrollo, economías ricas en recursos y casos africanos (Ajide, 2022; Avom et al., 2022; Destek et al., 2023; Kelly et al., 2024; Owjimehr \& Jamshidi, 2024; Valentine et al., 2024; Zhang et al., 2023), así como análisis de causalidad o no linealidad entre complejidad y rentas de recursos naturales (Olaniyi \& Odhiambo, 2025b; Ul-Durar et al., 2023). Frente a esos antecedentes, el diseño propuesto se centra en países dependientes de petróleo, gas y minería durante el período principal de estimación 1996--2022, estima especificaciones de datos de panel con efectos fijos para controlar por heterogeneidad no observable entre países y por shocks temporales comunes, y evalúa la sensibilidad de los resultados frente a tres umbrales de dependencia externa de recursos naturales ($\theta=20\%, 30\%, 40\%$). En conjunto, la contribución esperada es ofrecer evidencia comparada y organizada sobre las condiciones bajo las cuales las rentas extractivas se asocian con mayores o menores niveles de transformación estructural externa, sin sobredimensionar el alcance causal del diseño ni presentar como inéditos elementos que la literatura ya ha trabajado.

\section*{Referencias seleccionadas}

\begingroup
\small
\setlength{\parindent}{0pt}
\setlength{\parskip}{0.8em}

\hangindent=1.27cm\hangafter=1
Ajide, K. B. (2022). Is natural resource curse thesis an empirical regularity for economic complexity in Africa? \textit{Resources Policy, 76}, 1--11.

\hangindent=1.27cm\hangafter=1
Alvarado, R., Murshed, M., Cifuentes-Faura, J., Işık, C., Hossain, M. R., \& Tillaguango, B. (2023). Nexuses between rent of natural resources, economic complexity, and technological innovation: The roles of GDP, human capital and civil liberties. \textit{Resources Policy, 85}.

\hangindent=1.27cm\hangafter=1
Avom, D., Keneck-Massil, J., Njangang, H., \& Nvuh-Njoya, Y. (2022). Why are some resource-rich countries more sophisticated than others? The role of the regime type and political ideology. \textit{Resources Policy, 79}, 1--15.

\hangindent=1.27cm\hangafter=1
Avom, D., \& Ndoya, H. (2024). Does country stability spur economic complexity? Evidence from panel data. \textit{Applied Economics Letters, 31}(4), 323--329.

\hangindent=1.27cm\hangafter=1
Canh, N. P., Schinckus, C., \& Thanh, S. D. (2020). The natural resources rents: Is economic complexity a solution for resource curse? \textit{Resources Policy, 69}, 1--12.

\hangindent=1.27cm\hangafter=1
Destek, M. A., Hossain, M. R., Aydın, S., Shakib, M., \& Destek, G. (2023). Investigating the role of economic complexity in evading the resource curse. \textit{Resources Policy, 86}.

\hangindent=1.27cm\hangafter=1
Hoang, D. P., Chu, L. K., \& To, T. T. (2023). How do economic policy uncertainty, geopolitical risk, and natural resources rents affect economic complexity? Evidence from advanced and emerging market economies. \textit{Resources Policy, 85}, 1--15.

\hangindent=1.27cm\hangafter=1
Kazemzadeh, E., Fuinhas, J. A., Shirazi, M., Koengkan, M., \& Silva, N. (2023). Does economic complexity increase energy intensity? \textit{Energy Efficiency}, 1--22.

\hangindent=1.27cm\hangafter=1
Kelly, A. M., Ketu, I., Tchapchet Tchouto, J. E., \& Ndeffo, L. N. (2024). Investigating the link between exhaustion of natural resources and economic complexity in sub-Saharan Africa. \textit{African Development Review, 36}(3), 486--502.

\hangindent=1.27cm\hangafter=1
Li, Z., Doğan, B., Ghosh, S., Chen, W. M., \& Balsalobre Lorente, D. (2024). Economic complexity, natural resources and economic progress in the era of sustainable development: Findings in the context of resource deployment challenges. \textit{Resources Policy, 88}, 1--13.

\hangindent=1.27cm\hangafter=1
Mayer, S. S., Ibrid, A. A., Kalf, H. A. I., Altememy, H. A., Al-Muttar, M. Y. O., Hammoode, J. A., Farhan, M. A., \& Jalil, S. H. (2023). Is financial development having an impact on economic complexity? Empirical study of Iraq. \textit{Cuadernos de Economía, 46}(131), 135--144.

\hangindent=1.27cm\hangafter=1
Olaniyi, C. O., \& Odhiambo, N. M. (2025a). Does economic complexity provide antidotal pathways to evade the resource curse syndrome? A novel role for institutions. \textit{Sustainable Development, 33}, 3118--3142.

\hangindent=1.27cm\hangafter=1
Olaniyi, C. O., \& Odhiambo, N. M. (2025b). Exploring novelties in the causal relationship between economic complexity and natural resource rent: Empirical insights from Nigeria and South Africa. \textit{Journal of Open Innovation: Technology, Market, and Complexity, 11}(1).

\hangindent=1.27cm\hangafter=1
Olaniyi, C. O., \& Odhiambo, N. M. (2025c). Finding explanations for weak economic complexity in resource-rich African countries: Exploring the role of natural resource endowment and institutional quality. \textit{Resources Policy, 101}, 1--23.

\hangindent=1.27cm\hangafter=1
Owjimehr, S., \& Jamshidi, N. (2024). Why do natural resource-rich economies resist improving the economic complexity? \textit{Regional Science Policy and Practice, 16}(7).

\hangindent=1.27cm\hangafter=1
Sweet, C., \& Eterovic, D. (2019). Do patent rights matter? 40 years of innovation, complexity and productivity. \textit{World Development, 115}, 78--93.

\hangindent=1.27cm\hangafter=1
Ul-Durar, S., Arshed, N., Anwar, A., Sharif, A., \& Liu, W. (2023). How does economic complexity affect natural resource extraction in resource rich countries? \textit{Resources Policy, 86}, 1--15.

\hangindent=1.27cm\hangafter=1
Valentine, S. B., Motande, M. M. N. I., \& Salim Ahmed, V. M. (2024). Revisiting natural resources and economic complexity nexus: Does financial development matter in developing countries? \textit{Resources Policy, 93}.

\hangindent=1.27cm\hangafter=1
Valverde Carbonell, J., Pietrobelli, C., \& Menéndez de Medina, M. (2023). Critical minerals and countries' mining competitiveness: An estimate through economic complexity techniques. \textit{Maastricht Economic and Social Research Institute on Innovation and Technology (UNU-MERIT), United Nations University}, 1--21.

\hangindent=1.27cm\hangafter=1
Yalta, A. Y., \& Yalta, T. (2021). Determinants of economic complexity in MENA countries. \textit{Journal of Emerging Economies and Policy, 6}(1), 5--16.

\hangindent=1.27cm\hangafter=1
Zhang, Y., Qayyum, M., \& Yu, Y. (2023). Effects of resource abundance on economic complexity: Evidence from spatial panel model. \textit{Journal of Cleaner Production, 427}.

\endgroup

\end{document}
